\chapter{Géométrie dans l'espace (introduction)}
\textit{Jusqu'ici, tu as travaillé dans le plan. La géométrie de l'espace ajoute une
troisième dimension : les figures ne tiennent plus sur une feuille. Ce chapitre pose
le vocabulaire (points, droites, plans), les règles qui disent comment ces objets se
rencontrent, puis les outils pour calculer dans l'espace (coordonnées, vecteurs). On
reste \emph{introductif} : l'objectif est de \og voir dans l'espace\fg{} et de
raisonner proprement sur des positions relatives.}

\section{Points, droites et plans de l'espace}

\begin{definition}[Objets de base]
L'\textbf{espace} est l'ensemble de tous les \textbf{points}. On y distingue trois
types d'objets :
\begin{itemize}
  \item les \textbf{points}, notés $A,B,C,\dots$ ;
  \item les \textbf{droites}, notées $(AB)$, $(d)$, $\Delta,\dots$ ;
  \item les \textbf{plans}, notés $(ABC)$, $(\mathcal{P})$, $(\mathcal{Q}),\dots$
\end{itemize}
Un point peut \textbf{appartenir} ($\in$) ou non à une droite ou à un plan ; une
droite peut être \textbf{incluse} ($\subset$) ou non dans un plan.
\end{definition}

\begin{remarque}
On dessine l'espace en \textbf{perspective cavalière} : les figures du plan «~du
fond~» sont vues de face (vraies grandeurs), et la profondeur est représentée par des
fuyantes obliques. Les arêtes \emph{cachées} se tracent en pointillés.
\end{remarque}

\begin{aretenir}[Vocabulaire de position]
\begin{itemize}
  \item Des points sont \textbf{alignés} s'ils appartiennent à une même droite.
  \item Des points sont \textbf{coplanaires} s'ils appartiennent à un même plan.
  \item Des droites sont \textbf{coplanaires} s'il existe un plan les contenant toutes.
\end{itemize}
\end{aretenir}

\section{Règles d'incidence}

Les règles d'incidence disent \emph{quand} et \emph{comment} les objets se déterminent
ou se rencontrent. Elles sont admises (axiomes de la géométrie de l'espace).

\begin{propriete}[Déterminer une droite, déterminer un plan]
\begin{itemize}
  \item Par \textbf{deux points distincts} passe une \textbf{unique droite}.
  \item Par \textbf{trois points non alignés} passe un \textbf{unique plan}.
\end{itemize}
On détermine aussi un unique plan par : une droite et un point hors de cette droite ;
deux droites sécantes ; deux droites strictement parallèles.
\end{propriete}

\begin{propriete}[Droite incluse dans un plan]
Si deux points distincts $A$ et $B$ appartiennent à un plan $(\mathcal{P})$, alors
\textbf{toute la droite} $(AB)$ est incluse dans $(\mathcal{P})$.
\end{propriete}

\begin{theoreme}[Intersection de deux plans]
Si deux plans distincts ont un point commun, alors leur intersection est une
\textbf{droite} (la droite contient ce point, et tous les autres points communs).
\end{theoreme}

\begin{exemple}
Sur un cube $ABCDEFGH$, les points $A$, $B$ et $C$ ne sont pas alignés : ils
déterminent le plan de la face inférieure, noté $(ABC)$, qui contient aussi le point
$D$. La droite $(AB)$ est entièrement incluse dans ce plan.
\end{exemple}

\begin{illus}
\begin{tikzpicture}[scale=1.05,line join=round]
  \coordinate (A) at (0,0);   \coordinate (B) at (3,0);
  \coordinate (C) at (3.8,0.9); \coordinate (D) at (0.8,0.9);
  \coordinate (E) at (0,2.4);  \coordinate (F) at (3,2.4);
  \coordinate (G) at (3.8,3.3); \coordinate (H) at (0.8,3.3);
  % faces visibles
  \draw[thick,onbink] (A)--(B)--(F)--(E)--cycle;
  \draw[thick,onbink] (B)--(C)--(G)--(F)--cycle;
  \draw[thick,onbink] (E)--(F)--(G)--(H)--cycle;
  % arêtes cachées
  \draw[onbmuted,dashed] (A)--(D)--(C);
  \draw[onbmuted,dashed] (D)--(H);
  \foreach \p/\pos in {A/below left,B/below right,C/right,D/left,E/left,F/right,G/right,H/above left}
    \filldraw[onbo] (\p) circle (1.3pt) node[\pos,onbink,font=\footnotesize]{$\p$};
\end{tikzpicture}
\fig{Figure — un cube $ABCDEFGH$ en perspective cavalière : arêtes visibles en trait plein, arêtes cachées en pointillés.}
\end{illus}

\section{Positions relatives de deux droites}

\begin{definition}[Droites coplanaires / non coplanaires]
Deux droites de l'espace sont :
\begin{itemize}
  \item \textbf{coplanaires} s'il existe un plan les contenant toutes les deux ;
  \item \textbf{non coplanaires} sinon. On dit aussi qu'elles sont
        \textbf{gauches}.
\end{itemize}
\end{definition}

\begin{propriete}[Les trois cas]
Soit $(d)$ et $(d')$ deux droites de l'espace.
\begin{itemize}
  \item Si elles sont \textbf{coplanaires} : elles sont soit \textbf{sécantes}
        (un unique point commun), soit \textbf{parallèles} (strictement parallèles
        sans point commun, ou confondues).
  \item Si elles sont \textbf{non coplanaires} : elles n'ont \textbf{aucun point
        commun} et ne sont \textbf{pas} parallèles.
\end{itemize}
\end{propriete}

\begin{aretenir}
\textbf{Attention :} dans l'espace, deux droites sans point commun ne sont pas
forcément parallèles ! Elles peuvent être \textbf{non coplanaires}. Il ne suffit
pas de « ne pas se couper » pour être parallèles : il faut \emph{en plus} être
coplanaires.
\end{aretenir}

\begin{exemple}
Sur le cube ci-dessus : $(AB)$ et $(EF)$ sont parallèles ; $(AB)$ et $(BC)$ sont
sécantes (en $B$) ; $(AB)$ et $(CG)$ sont \textbf{non coplanaires} : aucun plan ne
les contient toutes les deux, et pourtant elles ne se coupent pas.
\end{exemple}

\section{Position relative d'une droite et d'un plan}

\begin{definition}[Les trois cas]
Soit une droite $(d)$ et un plan $(\mathcal{P})$.
\begin{itemize}
  \item $(d)$ est \textbf{incluse} dans $(\mathcal{P})$ : tous ses points sont dans
        $(\mathcal{P})$ ;
  \item $(d)$ est \textbf{parallèle} à $(\mathcal{P})$ : elle n'a aucun point commun
        avec $(\mathcal{P})$ ;
  \item $(d)$ est \textbf{sécante} à $(\mathcal{P})$ : elle a un \textbf{unique}
        point commun avec $(\mathcal{P})$.
\end{itemize}
\end{definition}

\begin{theoreme}[Critère de parallélisme droite/plan]
Une droite $(d)$ est parallèle à un plan $(\mathcal{P})$ \textbf{si et seulement si}
$(d)$ est parallèle à \textbf{une} droite contenue dans $(\mathcal{P})$.
\end{theoreme}

\begin{exemple}
La droite $(EH)$ du cube est parallèle au plan de la face $(ABC)$ : en effet $(EH)$
est parallèle à la droite $(AD)$, qui est, elle, contenue dans le plan $(ABC)$.
\end{exemple}

\section{Position relative de deux plans}

\begin{definition}[Les deux cas]
Soit deux plans $(\mathcal{P})$ et $(\mathcal{Q})$.
\begin{itemize}
  \item Ils sont \textbf{parallèles} s'ils n'ont aucun point commun (parallèles
        stricts) ou s'ils sont confondus ;
  \item ils sont \textbf{sécants} si leur intersection est une \textbf{droite}.
\end{itemize}
\end{definition}

\begin{propriete}[Critère de parallélisme de deux plans]
Si deux droites sécantes d'un plan $(\mathcal{P})$ sont respectivement parallèles à
deux droites d'un plan $(\mathcal{Q})$, alors les plans $(\mathcal{P})$ et
$(\mathcal{Q})$ sont parallèles.
\end{propriete}

\begin{remarque}
Sur le cube, les faces opposées sont parallèles : $(ABC)\parallel(EFG)$. Deux faces
adjacentes sont sécantes : leur intersection est l'arête commune (par exemple
$(ABC)\cap(BCG)=(BC)$).
\end{remarque}

\section{Parallélisme : le théorème du toit}

\begin{theoreme}[Théorème du toit]
Soit deux droites \textbf{parallèles} $(d)$ et $(d')$, l'une contenue dans un plan
$(\mathcal{P})$, l'autre dans un plan $(\mathcal{Q})$. Si $(\mathcal{P})$ et
$(\mathcal{Q})$ sont \textbf{sécants} selon une droite $\Delta$, alors $\Delta$ est
\textbf{parallèle} à $(d)$ et à $(d')$.
\end{theoreme}

\begin{illus}
\begin{tikzpicture}[scale=1.0,line join=round]
  % deux "pans de toit" sécants selon la ligne de faîte
  \coordinate (Pl) at (-2.2,0);  \coordinate (Pr) at (0.6,0);
  \coordinate (Ql) at (1.4,0);   \coordinate (Qr) at (4.2,0);
  \coordinate (Fa) at (-0.2,1.9); \coordinate (Fb) at (2.6,1.9);
  \fill[onbsoft] (Pl)--(Pr)--(Fb)--(Fa)--cycle;
  \fill[onbsoft] (Ql)--(Qr)--(Fb)--(Fa)--cycle;
  \draw[thick,onbink] (Pl)--(Pr)--(Fb)--(Fa)--cycle;
  \draw[thick,onbink] (Ql)--(Qr)--(Fb)--(Fa)--cycle;
  % droites parallèles dans chaque plan + faîte
  \draw[thick,onbo] (Pl)--(Pr) node[below,font=\footnotesize]{$(d)$};
  \draw[thick,onbo] (Qr)--(Ql) node[below,font=\footnotesize]{$(d')$};
  \draw[very thick,onbgreen] (Fa)--(Fb) node[midway,above,font=\footnotesize]{$\Delta$};
\end{tikzpicture}
\fig{Figure — théorème du toit : $(d)\parallel(d')$ entraîne $\Delta\parallel(d)\parallel(d')$, où $\Delta$ est l'arête (faîtage) commune aux deux plans.}
\end{illus}

\section{Orthogonalité : droite orthogonale à un plan}

\begin{definition}[Droite orthogonale à un plan]
Une droite $(d)$ est \textbf{orthogonale} à un plan $(\mathcal{P})$ si elle est
orthogonale à \textbf{toutes} les droites de $(\mathcal{P})$.
\end{definition}

\begin{theoreme}[Critère pratique]
Pour montrer qu'une droite $(d)$ est orthogonale à un plan $(\mathcal{P})$, il suffit
qu'elle soit orthogonale à \textbf{deux droites sécantes} de $(\mathcal{P})$.
\end{theoreme}

\begin{propriete}
Si une droite $(d)$ est orthogonale à un plan $(\mathcal{P})$, alors :
\begin{itemize}
  \item toute droite parallèle à $(d)$ est aussi orthogonale à $(\mathcal{P})$ ;
  \item tout plan parallèle à $(\mathcal{P})$ est aussi orthogonal à $(d)$.
\end{itemize}
\end{propriete}

\begin{exemple}
Sur le cube, l'arête $(AE)$ est orthogonale au plan $(ABC)$ : elle est orthogonale à
$(AB)$ et à $(AD)$, qui sont deux droites sécantes de ce plan. Elle est donc
orthogonale à \emph{toute} droite de la face $(ABC)$.
\end{exemple}

\section{Repérage dans l'espace}

\begin{definition}[Repère, coordonnées]
Un \textbf{repère} de l'espace est un quadruplet $(O\,;\,\vec{\imath},\vec{\jmath},
\vec{k})$ formé d'un point $O$ (l'origine) et de trois vecteurs non coplanaires.
Tout point $M$ de l'espace se repère alors par un unique triplet $(x\,;\,y\,;\,z)$,
ses \textbf{coordonnées}, tel que
\[ \overrightarrow{OM}=x\,\vec{\imath}+y\,\vec{\jmath}+z\,\vec{k}. \]
On dit que $x$ est l'\textbf{abscisse}, $y$ l'\textbf{ordonnée}, $z$ la \textbf{cote}.
\end{definition}

\begin{illus}
\begin{tikzpicture}[scale=1.1,line join=round]
  \coordinate (O) at (0,0);
  \draw[->,thick,onbink] (O)--(3,0) node[right,font=\footnotesize]{$x$};
  \draw[->,thick,onbink] (O)--(0,2.6) node[above,font=\footnotesize]{$z$};
  \draw[->,thick,onbink] (O)--(-1.5,-1.1) node[left,font=\footnotesize]{$y$};
  % point M(2;1;1.5) en perspective
  \coordinate (Mx) at (2,0);
  \coordinate (Mxy) at (1.1,-0.55);
  \coordinate (M) at (1.1,1.05);
  \draw[onbmuted,dashed] (O)--(Mx)--(Mxy)--(M);
  \draw[onbmuted,dashed] (Mxy)--(0,0);
  \filldraw[onbo] (M) circle (1.5pt) node[right,onbink,font=\footnotesize]{$M(x\,;\,y\,;\,z)$};
\end{tikzpicture}
\fig{Figure — repère $(O\,;\,\vec{\imath},\vec{\jmath},\vec{k})$ et coordonnées $(x\,;\,y\,;\,z)$ d'un point $M$.}
\end{illus}

\begin{propriete}[Milieu et distance (repère orthonormé)]
Dans un repère \textbf{orthonormé}, pour $A(x_A\,;\,y_A\,;\,z_A)$ et
$B(x_B\,;\,y_B\,;\,z_B)$ :
\begin{itemize}
  \item le \textbf{milieu} $I$ de $[AB]$ a pour coordonnées
        \[ I\!\left(\frac{x_A+x_B}{2}\,;\,\frac{y_A+y_B}{2}\,;\,\frac{z_A+z_B}{2}\right); \]
  \item la \textbf{distance} $AB$ vaut
        \[ AB=\sqrt{(x_B-x_A)^2+(y_B-y_A)^2+(z_B-z_A)^2}. \]
\end{itemize}
\end{propriete}

\begin{exemple}
Pour $A(1\,;\,0\,;\,2)$ et $B(3\,;\,4\,;\,2)$ : le milieu est $I(2\,;\,2\,;\,2)$ et
\[ AB=\sqrt{(3-1)^2+(4-0)^2+(2-2)^2}=\sqrt{4+16+0}=\sqrt{20}=2\sqrt{5}. \]
\end{exemple}

\section{Vecteurs de l'espace}

Les vecteurs du plan se prolongent sans changement à l'espace : on additionne, on
multiplie par un réel de la même façon. Seule nouveauté : il y a maintenant
\emph{trois} directions de base.

\begin{definition}[Colinéarité]
Deux vecteurs $\vec{u}$ et $\vec{v}$ sont \textbf{colinéaires} s'il existe un réel
$k$ tel que $\vec{v}=k\,\vec{u}$ (avec $\vec{u}\neq\vec{0}$). En coordonnées,
$\vec{u}(x\,;\,y\,;\,z)$ et $\vec{v}(x'\,;\,y'\,;\,z')$ sont colinéaires lorsque
$x'=kx,\ y'=ky,\ z'=kz$ pour un même $k$.
\end{definition}

\begin{propriete}
Trois points $A$, $B$, $C$ sont \textbf{alignés} si et seulement si les vecteurs
$\overrightarrow{AB}$ et $\overrightarrow{AC}$ sont colinéaires.
\end{propriete}

\begin{definition}[Vecteurs coplanaires]
Trois vecteurs $\vec{u}$, $\vec{v}$, $\vec{w}$ sont \textbf{coplanaires} si l'un
d'eux s'écrit comme combinaison des deux autres, par exemple
\[ \vec{w}=a\,\vec{u}+b\,\vec{v}\qquad (a,b\in\R). \]
Géométriquement, en plaçant leurs origines en un même point, leurs supports sont
contenus dans un même plan.
\end{definition}

\begin{remarque}
Quatre points $A$, $B$, $C$, $D$ sont coplanaires si et seulement si les vecteurs
$\overrightarrow{AB}$, $\overrightarrow{AC}$, $\overrightarrow{AD}$ sont coplanaires.
Cette idée sera approfondie en Terminale ; ici on se contente de la reconnaître.
\end{remarque}

\begin{exemple}
Soit $\vec{u}(1\,;\,0\,;\,1)$, $\vec{v}(0\,;\,1\,;\,1)$ et $\vec{w}(2\,;\,1\,;\,3)$.
On a $\vec{w}=2\,\vec{u}+1\,\vec{v}$ car $2(1\,;\,0\,;\,1)+(0\,;\,1\,;\,1)=
(2\,;\,1\,;\,3)$. Les trois vecteurs sont donc \textbf{coplanaires}.
\end{exemple}

\section{Sections planes simples}

\begin{definition}[Section plane]
La \textbf{section} d'un solide par un plan $(\mathcal{P})$ est l'ensemble des points
communs au solide et au plan. Pour un cube ou une pyramide, c'est un \textbf{polygone}.
\end{definition}

\begin{aretenir}[Construire une section]
Pour tracer la section d'un solide par un plan, on cherche, face par face, les
\textbf{segments} où le plan coupe chaque face :
\begin{itemize}
  \item deux points du plan situés sur une même face se relient par un segment ;
  \item on utilise le \textbf{parallélisme} : si le plan coupe deux faces parallèles,
        les deux traces obtenues sont parallèles.
\end{itemize}
La section est le polygone formé par tous ces segments mis bout à bout.
\end{aretenir}

\begin{exemple}
Le plan passant par les milieux $I$, $J$, $K$ de trois arêtes issues d'un même
sommet d'un cube coupe le cube selon un \textbf{triangle} $IJK$.
\end{exemple}

\begin{illus}
\begin{tikzpicture}[scale=1.05,line join=round]
  % pyramide à base carrée SABCD
  \coordinate (A) at (0,0);   \coordinate (B) at (3,0);
  \coordinate (C) at (3.8,0.9); \coordinate (D) at (0.8,0.9);
  \coordinate (S) at (1.9,3.0);
  \draw[onbmuted,dashed] (A)--(D)--(C);
  \draw[thick,onbink] (A)--(B)--(C);
  \draw[thick,onbink] (S)--(A); \draw[thick,onbink] (S)--(B); \draw[thick,onbink] (S)--(C);
  \draw[onbmuted,dashed] (S)--(D);
  % section parallèle à la base : milieux des arêtes latérales
  \coordinate (P) at ($(S)!0.5!(A)$);
  \coordinate (Q) at ($(S)!0.5!(B)$);
  \coordinate (R) at ($(S)!0.5!(C)$);
  \coordinate (T) at ($(S)!0.5!(D)$);
  \draw[thick,onbo] (P)--(Q)--(R);
  \draw[onbo,dashed] (R)--(T)--(P);
  \foreach \p/\pos in {S/above,A/below left,B/below right,C/right,D/left}
    \filldraw[onbink] (\p) circle (1.2pt) node[\pos,font=\footnotesize]{$\p$};
\end{tikzpicture}
\fig{Figure — pyramide $SABCD$ : un plan parallèle à la base la coupe selon un quadrilatère $PQRT$ semblable à la base.}
\end{illus}

% ════════════════════════════════════════════════════════════
\section{L'essentiel à retenir}

\begin{synthese}
\textbf{Détermination.} Deux points distincts $\Rightarrow$ une droite ; trois points
non alignés $\Rightarrow$ un plan. Si $A,B\in(\mathcal{P})$ alors $(AB)\subset(\mathcal{P})$.

\smallskip
\textbf{Deux droites.} Coplanaires : sécantes ou parallèles. Non coplanaires
(gauches) : ni point commun, ni parallèles.

\smallskip
\textbf{Droite / plan.} Incluse, parallèle, ou sécante (un point). Critère :
$(d)\parallel(\mathcal{P})\iff (d)$ parallèle à une droite de $(\mathcal{P})$.

\smallskip
\textbf{Deux plans.} Parallèles, ou sécants selon une droite.

\smallskip
\textbf{Théorème du toit.} $(d)\parallel(d')$, $(d)\subset(\mathcal{P})$,
$(d')\subset(\mathcal{Q})$, $(\mathcal{P})\cap(\mathcal{Q})=\Delta$ $\Rightarrow$
$\Delta\parallel(d)\parallel(d')$.

\smallskip
\textbf{Orthogonalité.} $(d)\perp(\mathcal{P})$ si $(d)$ orthogonale à deux droites
sécantes de $(\mathcal{P})$.

\smallskip
\textbf{Repère orthonormé.} $I=\big(\tfrac{x_A+x_B}{2};\tfrac{y_A+y_B}{2};
\tfrac{z_A+z_B}{2}\big)$ ; $AB=\sqrt{(x_B-x_A)^2+(y_B-y_A)^2+(z_B-z_A)^2}$.

\smallskip
\textbf{Vecteurs.} Colinéaires : $\vec{v}=k\vec{u}$ (alignement). Coplanaires :
$\vec{w}=a\vec{u}+b\vec{v}$.
\end{synthese}

% ════════════════════════════════════════════════════════════
\section{Méthodes \& savoir-faire}

\methode{Montrer que des points sont coplanaires ou alignés}
Pour l'alignement de $A,B,C$ : vérifier que $\overrightarrow{AB}$ et
$\overrightarrow{AC}$ sont colinéaires. Pour la coplanarité de $A,B,C,D$ : chercher
$a,b$ tels que $\overrightarrow{AD}=a\,\overrightarrow{AB}+b\,\overrightarrow{AC}$.
\begin{exemple}
$A(0;0;0)$, $B(1;1;0)$, $C(2;2;0)$. On a $\overrightarrow{AC}=(2;2;0)=2\,
\overrightarrow{AB}$ : $\overrightarrow{AB}$ et $\overrightarrow{AC}$ sont
colinéaires, donc $A$, $B$, $C$ sont alignés.
\end{exemple}

\methode{Reconnaître la position relative de deux droites}
Chercher d'abord un \emph{point commun}. S'il y en a un, elles sont sécantes. Sinon,
tester si elles sont \emph{coplanaires} (vecteurs directeurs colinéaires
$\Rightarrow$ parallèles ; sinon, vérifier l'existence d'un plan commun) : si aucun
plan ne les contient, elles sont \textbf{non coplanaires}.
\begin{exemple}
Sur un cube, $(AB)$ et $(CG)$ : pas de point commun, et leurs directions ne sont pas
parallèles ; on conclut qu'elles sont non coplanaires (gauches).
\end{exemple}

\methode{Montrer qu'une droite est parallèle à un plan}
Trouver, dans le plan $(\mathcal{P})$, une droite $(d')$ \textbf{parallèle} à la
droite $(d)$ étudiée. Souvent $(d')$ est une arête bien choisie du solide.
\begin{exemple}
$(EH)$ est parallèle au plan $(ABC)$ car $(EH)\parallel(AD)$ et $(AD)\subset(ABC)$.
\end{exemple}

\methode{Montrer qu'une droite est orthogonale à un plan}
Exhiber \textbf{deux droites sécantes} du plan, et montrer que la droite est
orthogonale à chacune. Conclure par le critère pratique.
\begin{exemple}
$(AE)\perp(ABC)$ : on vérifie $(AE)\perp(AB)$ et $(AE)\perp(AD)$, deux droites
sécantes de la face $(ABC)$.
\end{exemple}

\methode{Calculer une distance ou un milieu dans l'espace}
Placer les coordonnées des points, puis appliquer directement les formules du
repère orthonormé. Pour une longueur d'arête ou de diagonale, c'est la formule de
distance.
\begin{exemple}
Diagonale d'un cube d'arête $a$ avec $A(0;0;0)$ et $G(a;a;a)$ :
$AG=\sqrt{a^2+a^2+a^2}=a\sqrt{3}$.
\end{exemple}

\methode{Construire une section plane simple}
Marquer les points donnés sur les arêtes. Relier deux points d'une même face par un
segment. Utiliser le parallélisme : la trace sur deux faces parallèles est faite de
deux segments parallèles. Fermer le polygone.
\begin{exemple}
Un plan parallèle à la base d'une pyramide $SABCD$, passant par les milieux des
arêtes latérales, donne une section $PQRT$ qui est un \textbf{carré} (réduit) parallèle
à la base.
\end{exemple}
